%% Drafted from the mapped corpus by scripts/draft_topics.py.
% Model: gpt-5.6-luna; source characters: 24250.
\section{Процедурно програмиране}

\begin{defn}[Процедурно програмиране]
Процедурното програмиране е програмна парадигма, при която поведението на програмата се реализира чрез процедури, наричани още функции. Процедурите представляват самостоятелни части от програмата, които изпълняват определена задача и могат да бъдат извиквани многократно. Една процедура може да извиква друга процедура, поради което цялата програма се представя като поредица от стъпки, организирани в йерархия от извиквания.
\end{defn}

Процедурното програмиране се класифицира като императивно програмиране. При него програмата съдържа директни команди, които определят какво да бъде изпълнено и в какъв ред. Основни понятия са променливата, операторът, блокът, функцията и областта на действие.

\begin{defn}[Блок]
Блок е последователност от оператори, оградена с фигурни скоби $\{\}$.
Блокът се разглежда като една съставна конструкция и обикновено определя областта на действие на локално дефинираните в него променливи.
\end{defn}

\begin{defn}[Област на действие]
Областта на действие на даден идентификатор е частта от програмния текст, в която този идентификатор може да бъде използван.
\end{defn}

Понятията блок и област на действие са съществени за структурирането на програмата. Чрез тях се ограничава достъпът до данните и се отделят отделните части на алгоритъма.

\subsection{Принципи на структурното програмиране}

Основната идея на структурното програмиране е изграждането на ясна и разбираема структура на програмата. Програмата се разделя на отделни същности, всяка от които има ясно определена задача и граници. Пример за такава същност е функцията.

Основен метод за постигане на ясна структура е декомпозицията.

\begin{defn}[Декомпозиция]
Декомпозицията е разделяне на общата задача на по-малки подзадачи, всяка от които може да бъде реализирана чрез отделна процедура или функция.
\end{defn}

При декомпозиция сложната задача се представя чрез йерархия от по-прости задачи. Функциите, които реализират тези задачи, могат да бъдат извиквани от други функции. Така програмата се организира като множество от взаимосвързани, но логически обособени компоненти.

\begin{keydefn}[Модулност]
Модулността е принцип, според който функционалността на програмата се разпределя между отделни, взаимозаменяеми модули. Всеки модул трябва да съдържа всичко необходимо за изпълнението на един ясно определен аспект от програмата.
\end{keydefn}

Модулността има две основни цели:

\begin{itemize}
\item всяка част от програмата да има ясно определена отговорност;
\item отделните части да могат да бъдат използвани, проверявани и при необходимост заменяни независимо една от друга.
\end{itemize}

Добре структурираната процедурна програма се характеризира с последователност от ясно разграничени действия, функции с определено предназначение и ограничена област на действие на данните.

\section{Управление на изчислителния процес}

\begin{defn}[Изчислителен процес]
Изчислителният процес представлява изпълнение на поредица от стъпки в определен ред. Той има начална точка, последователност от междинни действия и крайна точка.
\end{defn}

В C++ изпълнението на програмата започва от функцията \texttt{main}. Функцията \texttt{main} може да извиква други функции. При нормално изпълнение програмата завършва, когато приключи изпълнението на \texttt{main}.

Една от специалните форми на дефиниция на \texttt{main} е:

\begin{verbatim}
int main(int argc, char* argv[])
\end{verbatim}

Параметърът \texttt{argc} указва броя на елементите в масива \texttt{argv}, а \texttt{argv} съдържа параметрите, подадени при стартиране на програмата от терминала. Възможна е и форма без параметри:

\begin{verbatim}
int main()
\end{verbatim}

Управлението на изчислителния процес се осъществява чрез управляващи конструкции. Те променят нормалния линеен ред на изпълнение, като позволяват:

\begin{itemize}
\item избор между различни действия;
\item многократно изпълнение на дадена последователност от оператори;
\item прекратяване на част от текущото изпълнение.
\end{itemize}

Основните управляващи конструкции са условните оператори и операторите за цикъл.

\section{Условни оператори}

Условният оператор позволява даден оператор или блок да бъде изпълнен само когато определено условие е вярно.

\subsection{Оператор \texttt{if}}

Общият синтаксис е:

\begin{verbatim}
if (условие_0)
    оператор_0
else if (условие_1)
    оператор_1
else
    оператор_else
\end{verbatim}

Частите \texttt{else if} могат да се повторят произволен брой пъти, а частта \texttt{else} е незадължителна и може да се срещне най-много веднъж.

Когато операторът се състои от повече от една команда, командите се обединяват в блок:

\begin{verbatim}
if (условие) {
    оператор_1
    оператор_2
}
\end{verbatim}

Изпълнението протича по следния начин:

\begin{enumerate}
\item Оценява се първото условие.
\item Ако то има стойност \texttt{true}, изпълнява се съответният оператор или блок.
\item След това останалите части на условния оператор се прескачат.
\item Ако условието е \texttt{false}, се проверява следващото условие в \texttt{else if}.
\item При първото условие със стойност \texttt{true} се изпълнява съответният оператор и проверката приключва.
\item Ако всички условия са \texttt{false} и съществува \texttt{else}, се изпълнява операторът след \texttt{else}.
\end{enumerate}

\begin{example}
Следната програма проверява знака на число:

\begin{verbatim}
int x;
cin >> x;

if (x > 0) {
    cout << "positive";
} else if (x < 0) {
    cout << "negative";
} else {
    cout << "zero";
}
\end{verbatim}

При изпълнението се извършва само един от трите блока.
\end{example}

Ако няма част \texttt{else} и всички условия са неверни, не се изпълнява нито един от операторите в конструкцията.

\subsection{Условна, или тернарна, операция}

Условната операция е триместна операция със синтаксис:

\[
\langle\text{булев израз}\rangle
\mathrel{?}
\langle\text{израз}_1\rangle
:
\langle\text{израз}_2\rangle
\]

Първо се оценява булевият израз. Ако резултатът е \texttt{true}, се оценява и връща резултатът на $\langle\text{израз}_1\rangle$. Ако резултатът е \texttt{false}, се оценява и връща резултатът на $\langle\text{израз}_2\rangle$.

Пример:

\begin{verbatim}
int x = (y < 2) ? y + 1 : y - 2;
\end{verbatim}

Ако $y=1$, на \texttt{x} се присвоява стойност $2$. Ако $y=2$, на \texttt{x} се присвоява стойност $0$.

Тернарната операция е подходяща, когато трябва да се избере една от две стойности. Тя връща резултат и затова може да се използва като част от по-голям израз.

\subsection{Оператор за многозначен избор \texttt{switch}}

Операторът \texttt{switch} служи за проверка на стойността на израз спрямо множество зададени стойности.

Синтаксисът има вида:

\begin{verbatim}
switch (израз) {
    case константен_израз_1:
        оператори_1
        break;

    case константен_израз_2:
        оператори_2
        break;

    default:
        оператори_default
}
\end{verbatim}

Оценяваната стойност на израза се сравнява със стойностите след етикетите \texttt{case}. Когато бъде намерено съвпадение, изпълнението започва от съответния \texttt{case}.

\begin{example}
\begin{verbatim}
int day;
cin >> day;

switch (day) {
    case 1:
        cout << "Monday";
        break;
    case 2:
        cout << "Tuesday";
        break;
    case 3:
        cout << "Wednesday";
        break;
    default:
        cout << "Unknown day";
}
\end{verbatim}
\end{example}

Ако не е намерен съответстващ \texttt{case} и съществува етикет \texttt{default}, се изпълняват операторите след \texttt{default}. Ако няма съвпадение и няма \texttt{default}, операторът приключва без изпълнение на case-блок.

След намиране на съответстващ \texttt{case} изпълнението продължава последователно надолу в блока на \texttt{switch}. Самото достигане на следващ етикет \texttt{case} не прекратява изпълнението. За прекратяване се използва операторът \texttt{break}.

\begin{supp}[Бележка]
Ако в съответстващия \texttt{case} липсва \texttt{break}, се осъществява последователно преминаване към следващите case-блокове. Това поведение е част от описаната семантика на оператора \texttt{switch}, а не грешка в синтаксиса.
\end{supp}

Изпълнението на \texttt{switch} приключва при едно от следните събития:

\begin{itemize}
\item достигнат е краят на блока на \texttt{switch};
\item изпълнен е оператор за управление на потока като \texttt{break} или \texttt{return};
\item нормалното изпълнение е прекъснато по друга причина.
\end{itemize}

Ползата от \texttt{switch} спрямо множество \texttt{if--else if--else} конструкции е, че управляващият израз се оценява веднъж и след това се сравнява с различните константни стойности.

\section{Оператори за цикъл}

Цикълът е конструкция за управление на изчислителния процес, чрез която даден оператор или блок се изпълнява многократно. Повторението продължава, докато условието има стойност \texttt{true}. Когато условието стане \texttt{false}, цикълът приключва.

В C++ основните оператори за цикъл са \texttt{for}, \texttt{while} и \texttt{do while}.

\subsection{Оператор \texttt{for}}

Синтаксисът е:

\begin{verbatim}
for (инициализация; условие; корекция)
    тяло
\end{verbatim}

При използване на блок:

\begin{verbatim}
for (инициализация; условие; корекция) {
    оператори
}
\end{verbatim}

Изпълнението протича в следния ред:

\begin{enumerate}
\item Изпълнява се инициализацията. Тя се изпълнява само веднъж при влизане в цикъла.
\item Оценява се условието.
\item Ако условието е \texttt{true}, изпълнява се тялото на цикъла.
\item След тялото се изпълнява корекцията.
\item Отново се оценява условието.
\item Процесът се повтаря, докато условието стане \texttt{false}.
\end{enumerate}

Най-често инициализацията служи за дефиниране и начално задаване на стойност на променлива, а корекцията увеличава или намалява тази променлива.

\begin{example}
\begin{verbatim}
for (int i = 0; i < 5; i++) {
    cout << i << " ";
}
\end{verbatim}

Изпълняват се стойностите $0$, $1$, $2$, $3$ и $4$. След като \texttt{i} стане $5$, условието \texttt{i < 5} е невярно и цикълът приключва.
\end{example}

Променлива, дефинирана в инициализационната част на \texttt{for}, има област на действие в цикъла, тоест може да бъде използвана до края на блока на цикъла.

\subsection{Оператор \texttt{while}}

Синтаксисът е:

\begin{verbatim}
while (условие)
    тяло
\end{verbatim}

или:

\begin{verbatim}
while (условие) {
    оператори
}
\end{verbatim}

При изпълнение първо се оценява условието. Ако то е \texttt{true}, се изпълнява тялото и след това отново се проверява условието. Ако условието е \texttt{false}, тялото не се изпълнява и цикълът приключва.

\begin{example}
\begin{verbatim}
int i = 0;

while (i < 5) {
    cout << i << " ";
    i++;
}
\end{verbatim}
\end{example}

При \texttt{while} е необходимо тялото на цикъла или друга част от изпълнението да води до промяна, която в даден момент да направи условието невярно. В противен случай повторението може да продължи без край.

\subsection{Оператор \texttt{do while}}

Операторът \texttt{do while} е сходен с \texttt{while}, но проверката на условието се извършва след изпълнението на тялото.

Синтаксисът е:

\begin{verbatim}
do {
    оператори
} while (условие);
\end{verbatim}

Редът на изпълнение е:

\begin{enumerate}
\item изпълнява се тялото на цикъла;
\item оценява се условието;
\item ако условието е \texttt{true}, тялото се изпълнява отново;
\item ако условието е \texttt{false}, цикълът приключва.
\end{enumerate}

Следователно тялото на \texttt{do while} се изпълнява поне веднъж, преди да бъде направена първата проверка.

\begin{example}
\begin{verbatim}
int x;

do {
    cin >> x;
} while (x < 0);
\end{verbatim}

Тялото се изпълнява, след което се проверява дали въведеното число е отрицателно.
\end{example}

\section{Променливи и присвояване}

\begin{defn}[Променлива]
Променливата е именувана област в паметта, свързана с име, тип, адрес и стойност.
\end{defn}

Името на променливата е идентификатор, чрез който тя се използва в програмния текст. Адресът указва мястото в паметта, отредено за променливата. Типът определя как компилаторът интерпретира съдържанието на това място и какви операции могат да се извършват със стойността. Стойността е текущото съдържание на променливата.

\subsection{Дефиниране и инициализиране}

Общият синтаксис за дефиниране на променливи е:

\[
\langle\text{тип}\rangle\ \langle\text{идентификатор}\rangle
[\mathrel{=}\langle\text{израз}\rangle]
\{
,\ \langle\text{идентификатор}\rangle
[\mathrel{=}\langle\text{израз}\rangle]
\}\ ;
\]

Например:

\begin{verbatim}
int a;
double x = 3.5;
int p = 1, q = 2;
\end{verbatim}

Дефинирането създава променлива от даден тип. Променливата е инициализирана, когато при създаването й е зададена начална стойност.

В C++ могат да се използват например следните форми:

\begin{verbatim}
int x = 5;
int y(5);
int z{5};
\end{verbatim}

При използване на присвояване с \texttt{=} се използва копираща инициализация. При формите с кръгли или фигурни скоби началната стойност се задава непосредствено при инициализацията.

\begin{supp}[Бележка]
Източниците описват неинициализираната локална променлива като съдържаща непредвидима стойност. Практическото следствие е, че такава променлива не трябва да се използва, преди да й бъде зададена определена стойност.
\end{supp}

\subsection{Оператор за присвояване}

Синтаксисът на присвояването е:

\[
\langle\text{идентификатор}\rangle
=
\langle\text{израз}\rangle;
\]

Операторът оценява израза отдясно и записва получената стойност в променливата отляво.

\begin{verbatim}
int x;
x = 7;
x = x + 1;
\end{verbatim}

Резултатът от присвояването е стойността, записана в променливата. Поради това присвояванията могат да се влагат едно в друго.

\begin{defn}[lvalue и rvalue]
$lvalue$ е място в паметта, което има стойност и може да бъде променяно. Обикновена променлива е пример за $lvalue$.

$rvalue$ е стойност, която не представлява самостоятелно променяемо място в паметта. Константа, литералът и резултатът от пресмятане са примери за $rvalue$.
\end{defn}

При израза

\[
a=b=c=2
\]

присвояването се извършва отдясно наляво:

\[
a=(b=(c=2)).
\]

Следователно първо на \texttt{c} се присвоява $2$, полученият резултат се присвоява на \texttt{b}, а след това и на \texttt{a}.

\subsection{Типове на променливите}

Типът указва на компилатора как да интерпретира съдържанието на паметта и носи семантична информация. Променливите от един и същи тип имат сходни характеристики, включително размер и допустими операции.

Типовете могат да бъдат разглеждани като:

\begin{itemize}
\item скаларни типове;
\item съставни типове.
\end{itemize}

Към скаларните типове спадат булевият, целочисленият, символният, изброимият тип, типовете за числа с плаваща запетая, указателите и референциите.

Към съставните типове спадат масивите, структурите и класовете.

\subsection{Локални и глобални променливи}

Областта на действие определя къде в програмния текст променливата може да бъде достъпвана.

\begin{defn}[Локална променлива]
Локална е променлива, дефинирана в блок или във функция. Областта й на действие започва от мястото на дефинирането и продължава до края на блока, в който е дефинирана.
\end{defn}

\begin{example}
\begin{verbatim}
int main() {
    int x = 1;

    if (x > 0) {
        int y = 2;
        cout << y;
    }

    // y вече не е достъпна тук
}
\end{verbatim}
\end{example}

\begin{defn}[Глобална променлива]
Глобална е променлива, дефинирана извън тялото на функция. Областта й на действие е от мястото на дефинирането до края на файла.
\end{defn}

Локалната област на действие ограничава използването на променливата до частта от програмата, в която тя е необходима. Глобалната променлива може да бъде достъпвана от по-голяма част от програмния файл.

\section{Функции и процедури}

\begin{defn}[Функция]
Функцията е самостоятелен фрагмент от програмата, съдържащ поредица от команди, предназначени за изпълнение на определена задача и позволяващи многократно използване.
\end{defn}

Функцията извършва пресмятане и връща резултат. Процедурата изпълнява поредица от оператори, без задължително да връща резултат. В C++ и двете понятия се реализират чрез функции. Функция, която не връща резултат, има резултатен тип \texttt{void}.

\subsection{Дефиниране на функция}

Общият синтаксис е:

\begin{verbatim}
тип_резултат име(формални_параметри) {
    тяло
}
\end{verbatim}

Например:

\begin{verbatim}
int square(int x) {
    return x * x;
}
\end{verbatim}

В дефиницията се съдържат:

\begin{itemize}
\item резултатен тип;
\item име на функцията;
\item списък от формални параметри;
\item тяло на функцията.
\end{itemize}

Резултатният тип може да бъде \texttt{void} или друг допустим тип, но в съвременния C++ трябва да бъде зададен изрично. Пропуснат резултатен тип не означава подразбиращ се \texttt{int}; такава стара конвенция е съществувала в ранни версии на C, но не е правило на C++.

Функция без параметри може да бъде представена с празен списък или с \texttt{void}:

\begin{verbatim}
void printMessage() {
    cout << "Hello";
}
\end{verbatim}

Формалните параметри имат вида:

\[
\langle\text{тип}\rangle\ \langle\text{идентификатор}\rangle
\{
,\ \langle\text{тип}\rangle\ \langle\text{идентификатор}\rangle
\}.
\]

Например:

\begin{verbatim}
int add(int a, int b) {
    return a + b;
}
\end{verbatim}

\subsection{Декларация и извикване}

Функцията може да бъде само декларирана, без да се задава тялото й. Декларацията представлява обещание, че дефиниция на функцията ще бъде предоставена.

Примерна декларация е:

\begin{verbatim}
int add(int a, int b);
\end{verbatim}

Извикването на функция се записва чрез нейното име и фактическите параметри:

\begin{verbatim}
име(фактически_параметри)
\end{verbatim}

Например:

\begin{verbatim}
int result = add(2, 3);
\end{verbatim}

Фактическите параметри, наричани още аргументи, са изрази:

\[
\langle\text{израз}\rangle
\{
,\ \langle\text{израз}\rangle
\}.
\]

При извикването типът на всеки фактически параметър се съпоставя с типа на съответния формален параметър. Ако е необходимо, се извършва преобразуване на типовете.

\subsection{Връщане на резултат}

Операторът за връщане има вида:

\begin{verbatim}
return израз;
\end{verbatim}

Той връща стойността на израза като резултат от извикването на функцията и незабавно прекратява изпълнението на функцията.

\begin{example}
\begin{verbatim}
int maximum(int a, int b) {
    if (a > b) {
        return a;
    }
    return b;
}
\end{verbatim}
\end{example}

Типът на израза след \texttt{return} се съпоставя с резултатния тип на функцията. Ако е необходимо, се извършва преобразуване.

Функция с резултатен тип \texttt{void} може да изпълнява оператори, без да връща стойност. В този случай работата й приключва при достигане на края на тялото или чрез \texttt{return} без резултат.

\subsection{Стекова рамка при извикване}

При извикване на функция се заделя нова стекова рамка. В нея се пазят фактическите параметри, адресът за връщане и локалните променливи на функцията. Рамката се намира на върха на програмния стек.

Следователно всяко извикване има собствен контекст. Локалните променливи и параметрите на едно извикване са свързани с неговата стекова рамка. При приключване на функцията изпълнението се връща към адреса, от който тя е била извикана.

\subsection{Предаване на параметри}

\subsubsection{Предаване по стойност}

При предаване по стойност първо се пресмята стойността на фактическия параметър. В стековата рамка на функцията се създава копие на тази стойност.

\begin{verbatim}
void change(int x) {
    x = 10;
}

int a = 3;
change(a);
\end{verbatim}

Промяната на \texttt{x} остава локална за функцията. Стойността на \texttt{a} не се променя. При завършване на функцията копието и направените върху него промени изчезват.

\subsubsection{Предаване чрез референция}

Предаването чрез референция се използва, когато промените във формалния параметър трябва да се отразят върху фактическия параметър.

Синтаксисът на параметър-референция е:

\begin{verbatim}
тип& идентификатор
\end{verbatim}

Пример:

\begin{verbatim}
int add5(int& x) {
    x += 5;
    return x;
}

int a = 3;
cout << add5(a) << " " << a;
\end{verbatim}

Резултатът е:

\begin{verbatim}
8 8
\end{verbatim}

В този случай параметърът \texttt{x} се свързва с променливата \texttt{a}. Затова промяната на \texttt{x} променя и \texttt{a}. При предаване чрез референция фактическият параметър трябва да бъде $lvalue$, тоест променяемо място в паметта.

\subsubsection{Съответствие на типовете}

По време на компилация се проверява съответствието между типовете на формалните и фактическите параметри. Проверява се и съответствието между типа на връщания израз и резултатния тип на функцията. Когато е необходимо, се извършва преобразуване на типовете.

\begin{keydefn}[Модел на нормално извикване и връщане]
При извикване на функция се оценяват фактическите параметри, съпоставят се със съответните формални параметри и се създава контекст за изпълнение на функцията. При нормално завършване функцията връща управлението към мястото на извикването. При функция, връщаща стойност, резултатът се задава чрез \texttt{return}; достигането на края е допустимо за \texttt{void} функции и за \texttt{main}. Описанието не определя реда на оценяване на отделните аргументи.
\end{keydefn}

Това твърдение обобщава връзката между процедурната декомпозиция, предаването на данни и управлението на изчислителния процес.

\section{Символни низове}

\begin{defn}[Символен низ]
Символният низ е последователност от символи. В C++ той се представя като масив от елементи от тип \texttt{char}, след последния символ на който се записва терминиращият символ \texttt{'\textbackslash 0'}.
\end{defn}

Терминиращият символ има код $0$ в ASCII таблицата. Той показва къде завършва низът, независимо от размера на масива, в който е записан.

Пример:

\begin{verbatim}
char word[] = {'H', 'e', 'l', 'l', 'o', '\0'};
\end{verbatim}

Същият низ може да бъде зададен чрез литерал:

\begin{verbatim}
char word[] = "Hello";
\end{verbatim}

При дефиниране на по-голям буфер:

\begin{verbatim}
char word1[100] = "Hello";
\end{verbatim}

се заделя място за $100$ символа. Низът заема символите на \texttt{Hello} и терминиращия символ, а останалото място остава част от масива. Максималната дължина на низа в такъв буфер е $99$ символа, тъй като един елемент трябва да бъде запазен за \texttt{'\textbackslash 0'}.

\subsection{Основни операции със символни низове}

Въвеждането чрез \texttt{cin >>} чете до разделител, например интервал, табулация или нов ред.

\begin{verbatim}
char word[100];
cin >> word;
\end{verbatim}

Функцията \texttt{getline} въвежда символи до нов ред, но трябва да й бъде зададен размер на буфера:

\begin{verbatim}
cin.getline(word, 100);
\end{verbatim}

В този случай се въвеждат най-много $99$ символа, така че да остане място за терминиращия символ.

Извеждането се извършва чрез \texttt{cout}:

\begin{verbatim}
cout << word;
\end{verbatim}

До отделен символ се достига чрез индексиране:

\begin{verbatim}
char word[] = "Hello";
cout << word[1];
\end{verbatim}

Резултатът е символът \texttt{'e'}, защото индексирането започва от $0$.

\subsection{Функции от библиотеката \texttt{cstring}}

За работа със символни низове се използват функции от заглавния файл \texttt{cstring}.

\begin{itemize}
\item \texttt{strlen(низ)} връща дължината на низа, тоест броя символи до първото срещане на \texttt{'\textbackslash 0'};
\item \texttt{strcpy(буфер, низ)} прехвърля символите на низа в буфера и връща буфера;
\item \texttt{strcmp(низ1, низ2)} сравнява два низа лексикографски;
\item \texttt{strcat(низ1, низ2)} добавя символите на \texttt{низ2} в края на \texttt{низ1};
\item \texttt{strchr(низ, символ)} търси първото срещане на символ и връща адрес към него или нулев указател;
\item \texttt{strstr(низ, подниз)} търси първото срещане на подниз и връща адрес към него или нулев указател.
\end{itemize}

Функцията \texttt{strcmp} връща:

\begin{itemize}
\item $0$, ако низовете съвпадат;
\item число, по-малко от $0$, ако \texttt{низ1} е преди \texttt{низ2};
\item число, по-голямо от $0$, ако \texttt{низ1} е след \texttt{низ2}.
\end{itemize}

При \texttt{strcat} програмистът трябва да осигури достатъчно място в първия буфер за символите на двата низа и за новия терминиращ символ.

Примерна реализация на \texttt{strlen} е:

\begin{verbatim}
int my_strlen(const char* str) {
    int length = 0;

    while (str[length] != '\0') {
        ++length;
    }

    return length;
}
\end{verbatim}

Алгоритъмът започва с дължина $0$ и последователно проверява символите на низа. При всеки символ, различен от \texttt{'\textbackslash 0'}, дължината се увеличава с единица. Когато бъде достигнат терминиращият символ, текущата дължина се връща.

\section{Изпитен фокус}

При устен отговор трябва да могат да бъдат възпроизведени:

\begin{itemize}
\item определението за процедурно програмиране и връзката му с императивното програмиране;
\item ролята на блоковете и областта на действие;
\item принципите на декомпозицията и модулността;
\item началото и краят на изпълнението чрез функцията \texttt{main};
\item синтаксисът и семантиката на \texttt{if}, \texttt{else if}, \texttt{else}, тернарния оператор и \texttt{switch};
\item поведението на \texttt{switch} при липса на \texttt{break};
\item синтаксисът и редът на изпълнение при \texttt{for}, \texttt{while} и \texttt{do while};
\item понятието променлива, нейните име, тип, адрес и стойност;
\item дефиниране, инициализиране и присвояване;
\item разликата между $lvalue$ и $rvalue$;
\item разликата между локална и глобална променлива;
\item видовете скаларни и съставни типове;
\item структурата на функцията, декларацията, дефиницията, извикването и \texttt{return};
\item предаването на параметри по стойност и чрез референция;
\item съпоставянето на типовете на формалните и фактическите параметри;
\item стековата рамка при извикване на функция;
\item представянето на символен низ чрез масив от \texttt{char} и терминиращ символ;
\item основните операции и функциите от \texttt{cstring};
\item алгоритъма за реализация на \texttt{strlen}.
\end{itemize}
